% ============================================================
% 第 17 章 同步发电机的基本理论（教材第五篇，印刷页 183--203）
% 节名：第一节 同步发电机的运行方式；第二节 隐极同步发电机的分析方法；
%       第三节 凸极同步发电机的分析方法；第四节 同步发电机功率和转矩方程；第五节 同步发电机的运行特性和电抗测定
% ============================================================
\chapter{同步发电机的基本理论}

同步发电机与变压器、感应电机的根本区别在于：它的主磁场由\textbf{转子上的直流励磁绕组}独立建立，而定子（电枢）三相电流产生的磁动势会反过来影响主磁场——这种现象称为\textbf{电枢反应}。整章的逻辑是：先确定电枢反应的性质（去磁、助磁还是交磁），再把电枢反应"折算"成电抗，于是就能像变压器那样写出电压方程、画相量图、算功率。

\section{同步发电机的运行方式}

同步发电机转子由原动机拖动，以同步转速 $n_1=60f/p$ 旋转。转子励磁绕组通入直流电流 $I_f$，建立\textbf{励磁磁动势} $F_f$ 与主磁场；主磁场切割定子三相对称绕组，感应三相对称电动势。定子接上负载后，三相电流产生\textbf{电枢磁动势} $F_a$，它与 $F_f$ 共同建立气隙磁场。

稳态对称运行时，只研究一相（通常取 $A$ 相），并用相量表示。各量正方向按发电机惯例规定：电枢电流 $\dot I$ 的正方向与感应电动势 $\dot E$ 相同（由中性点流向出线端），即电机\textbf{输出}电流为正。

\section{隐极同步发电机的分析方法}

\subsection{电枢反应的性质}
电枢磁动势 $F_a$ 对主磁场的作用称为\textbf{电枢反应}，其性质取决于负载的性质（即 $\dot I$ 与 $\dot E_0$ 的相位关系）。设 $\psi$ 为 $\dot E_0$ 与 $\dot I$ 的夹角（内功率因数角）：
\begin{itemize}
  \item $\psi=0$（纯电阻负载，$\dot I$ 与 $\dot E_0$ 同相）：$F_a$ 的轴线与主磁极轴线相差 $90^\circ$ 电角度，称为\textbf{交轴（交磁）电枢反应}，使气隙磁场轴线扭斜，磁场半边增强、半边削弱，不改变总磁通量；
  \item $\psi=90^\circ$（纯电感负载，$\dot I$ 滞后 $\dot E_0$ $90^\circ$）：$F_a$ 的轴线与主磁极轴线\textbf{相反}，称为\textbf{直轴去磁电枢反应}，使主磁通削弱、端电压下降（所以感性负载需要更大的励磁电流）；
  \item $\psi=-90^\circ$（纯电容负载，$\dot I$ 超前 $\dot E_0$ $90^\circ$）：$F_a$ 轴线与主磁极轴线\textbf{相同}，称为\textbf{直轴助磁电枢反应}，使主磁通增强、端电压升高。
\end{itemize}

\begin{center}
\begin{tabular}{@{}llll@{}}
\toprule
负载性质 & $\psi$ & 电枢反应性质 & 对主磁通/端电压的影响 \\
\midrule
纯电阻 & $0$ & 交轴（交磁） & 磁场扭斜，总量基本不变 \\
纯电感（感性） & $90^\circ$ & 直轴去磁 & 削弱主磁通，端电压下降 \\
纯电容（容性） & $-90^\circ$ & 直轴助磁 & 增强主磁通，端电压升高 \\
\bottomrule
\end{tabular}
\end{center}

\subsection{电枢反应电抗与同步电抗}
把电枢反应磁动势 $F_a$ 的作用等效为一个电抗压降：$F_a$ 产生的磁通在电枢绕组中感应电动势 $\dot E_a$，且 $\dot E_a=-\mathrm{j}\dot I X_a$（$X_a$ 为\textbf{电枢反应电抗}）。再加上漏磁通产生的漏抗压降 $-\mathrm{j}\dot I X_\sigma$，两者合并为
\begin{equation}
X_t=X_a+X_\sigma ,
\end{equation}
称为\textbf{同步电抗}。

\keyformula{X_t=X_a+X_\sigma}

$X_t$ 综合反映了电枢反应与漏磁两种效应。\textbf{电枢反应是同步电机区别于其他电机的核心}：它把"磁场互相作用"这个物理过程，压缩成一个电抗压降，从而使同步电机可以用电路方法来分析。这是"路的分析方法"最漂亮的一次应用。

\subsection{电压方程与相量图}
计及电枢电阻 $R_a$，隐极同步发电机的电压方程为
\begin{equation}
\dot E_0=\dot U+\dot I\left(R_a+\mathrm{j}X_t\right).
\end{equation}

\keyformula{\dot E_0=\dot U+\dot I\left(R_a+\mathrm{j}X_t\right)}

通常 $R_a\ll X_t$，可忽略，得
\begin{equation}
\dot E_0\approx\dot U+\mathrm{j}\dot I X_t .
\end{equation}

\subsubsection{相量图画法}
以端电压 $\dot U$ 为参考，按负载功率因数角 $\varphi$ 画出 $\dot I$；在 $\dot U$ 上加 $\dot I R_a$（平行 $\dot I$）与 $\mathrm{j}\dot I X_t$（超前 $\dot I$ $90^\circ$），其相量和即 $\dot E_0$；$\dot E_0$ 与 $\dot U$ 的夹角即\textbf{功率角}（功角）$\delta$，$\dot E_0$ 与 $\dot I$ 的夹角为 $\psi$，且有
\begin{equation}
\psi=\delta+\varphi .
\end{equation}

\section{凸极同步发电机的分析方法}

凸极电机气隙不均匀（极下小、极间大），同样大小的电枢磁动势作用在不同位置时产生的磁通不同，因此不能用一个同步电抗描述。工程上采用\textbf{双反应理论}：把电枢磁动势分解为两个分量，分别沿\textbf{直轴}（磁极轴线，$d$ 轴）与\textbf{交轴}（极间轴线，$q$ 轴），各自用相应的电抗计算。

\subsection{双反应理论}
把电枢电流分解为
\begin{equation}
\dot I=\dot I_d+\dot I_q,\qquad
I_d=I\sin\psi,\qquad I_q=I\cos\psi ,
\end{equation}
其中 $\dot I_d$ 产生直轴电枢反应（$X_{ad}$），$\dot I_q$ 产生交轴电枢反应（$X_{aq}$）。相应的同步电抗为
\begin{equation}
X_d=X_{ad}+X_\sigma,\qquad X_q=X_{aq}+X_\sigma .
\end{equation}
对凸极机通常 $X_d>X_q$（直轴气隙小、磁阻小，故电抗大）。

\subsection{电压方程}
\begin{equation}
\dot E_0=\dot U+\dot I R_a+\mathrm{j}\dot I_dX_d+\mathrm{j}\dot I_qX_q .
\end{equation}

\keyformula{\dot E_0=\dot U+\dot I R_a+\mathrm{j}\dot I_dX_d+\mathrm{j}\dot I_qX_q}

忽略 $R_a$ 时
\begin{equation}
\dot E_0\approx\dot U+\mathrm{j}\dot I_dX_d+\mathrm{j}\dot I_qX_q .
\end{equation}

\subsubsection{相量图的特点}
凸极机的相量图中，$\dot E_0$ 的位置不能像隐极机那样直接由 $\dot U+\mathrm{j}\dot I X_t$ 一次求出，因为 $X_d\ne X_q$、$I_d$、$I_q$ 又依赖 $\psi$。实用的作图法是：先由 $\dot U$ 与 $\dot I$ 求出 $\dot E_Q=\dot U+\dot I R_a+\mathrm{j}\dot I X_q$，其方向即 $q$ 轴方向（$\dot E_0$ 的方向），据此分解出 $I_d$、$I_q$，再求 $\dot E_0$。这一"辅助电动势 $\dot E_Q$"的技巧是解凸极机题目的关键。

\begin{center}
\begin{tabular}{@{}p{2.2cm}p{3.4cm}p{3.6cm}p{3.2cm}@{}}
\toprule
项目 & 隐极机 & 凸极机 & 说明 \\
\midrule
气隙 & 均匀 & 不均匀 & 决定能否用单一电抗 \\
分析方法 & 同步电抗 $X_t$ & 双反应理论（$X_d$、$X_q$） & 分解 $I_d$、$I_q$ \\
电压方程 & $\dot E_0=\dot U+\dot I(R_a+\mathrm{j}X_t)$ & $\dot E_0=\dot U+\dot IR_a+\mathrm{j}\dot I_dX_d+\mathrm{j}\dot I_qX_q$ & — \\
电抗关系 & $X_t=X_a+X_\sigma$ & $X_d=X_{ad}+X_\sigma$，$X_q=X_{aq}+X_\sigma$ & 通常 $X_d>X_q$ \\
相量图 & 一次求出 $\dot E_0$ & 需先求 $\dot E_Q$ 定 $q$ 轴 & 凸极机作图较繁 \\
\bottomrule
\end{tabular}
\end{center}

\corepoint{隐极与凸极的分界只有一句话：\textbf{气隙是否均匀}。均匀 $\Rightarrow$ 一个 $X_t$ 够了；不均匀 $\Rightarrow$ 必须分解 $d$、$q$ 两轴。考试若给了 $X_d$、$X_q$ 两个电抗，一定是凸极机；只给一个 $X_t$（或 $X_s$），一定是隐极机——先看参数再动笔，可省大量时间。}

\section{同步发电机功率和转矩方程}

\subsection{电磁功率}
发电机输出的有功功率为
\begin{equation}
P_2=mUI\cos\varphi ,
\end{equation}
加上电枢铜耗 $p_{Cu}=mI^2R_a$ 得电磁功率 $P_{em}=P_2+p_{Cu}$。忽略 $R_a$ 时 $P_{em}\approx P_2$。

\subsubsection{隐极机的功角特性}
由相量图的几何关系可得
\begin{equation}
P_{em}=\frac{mE_0U}{X_t}\sin\delta .
\end{equation}

\keyformula{P_{em}=\frac{mE_0U}{X_t}\sin\delta}

\subsubsection{凸极机的功角特性}
\begin{equation}
P_{em}=\frac{mE_0U}{X_d}\sin\delta+\frac{mU^{2}}{2}\left(\frac{1}{X_q}-\frac{1}{X_d}\right)\sin2\delta .
\end{equation}

\keyformula{P_{em}=\frac{mE_0U}{X_d}\sin\delta+\frac{mU^{2}}{2}\left(\frac1{X_q}-\frac1{X_d}\right)\sin2\delta}

两项的物理意义：
\begin{itemize}
  \item 第一项（$\sin\delta$）是\textbf{励磁电磁功率}，由励磁磁动势与电枢磁动势的相互作用产生，与 $E_0$（即励磁电流）成正比，是发电机传递有功的主要成分；
  \item 第二项（$\sin2\delta$）是\textbf{磁阻功率}（凸极效应功率），由 $d$、$q$ 轴磁阻不等产生，与励磁无关，即使 $E_0=0$ 也存在——这是\textbf{磁阻同步电动机}（无励磁转子）能工作的原理。
\end{itemize}

\subsection{转矩方程}
\begin{equation}
T=\frac{P_{em}}{\Omega_1},\qquad \Omega_1=\frac{2\pi n_1}{60}=\frac{2\pi f}{p}.
\end{equation}
稳态时 $T=T_1=T_2+T_0$（原动机拖动转矩 $=$ 电磁转矩 $=$ 输出转矩 $+$ 空载制动转矩）。

\section{同步发电机的运行特性和电抗测定}

\subsection{空载特性}
$n=n_1$、$I=0$ 时，端电压（空载电动势）$E_0$ 与励磁电流 $I_f$ 的关系 $E_0=f(I_f)$。由于铁芯饱和，曲线起始段近似直线（气隙线），随后弯曲并趋于饱和。空载特性本质上就是电机的磁化曲线。

\subsection{短路特性}
$n=n_1$、三相稳态短路（$U=0$）时，短路电流 $I_k$ 与励磁电流 $I_f$ 的关系 $I_k=f(I_f)$。由于短路时电枢反应为纯去磁、磁路不饱和，故\textbf{短路特性是一条直线}。

\subsection{外特性与调整特性}
\begin{itemize}
  \item \textbf{外特性}：$n=n_1$、$I_f=$ 常数、$\cos\varphi=$ 常数时，$U=f(I)$。感性负载时随 $I$ 增大 $U$ 下降；容性负载时可能上升。
  \item \textbf{调整特性}：$n=n_1$、$U=U_N$、$\cos\varphi=$ 常数时，$I_f=f(I)$。为维持端电压不变，感性负载需增加励磁。
\end{itemize}

\subsection{同步电抗与短路比}
由空载特性与短路特性可求同步电抗：在短路特性上取额定电流对应的励磁电流 $I_{fk}$，再在空载特性上查同一励磁电流下的空载电动势 $E_0$，则
\begin{equation}
X_d=\frac{E_0}{I_k}\quad\left(I_k=I_N\right).
\end{equation}
\textbf{短路比}定义为空载时产生额定电压所需的励磁电流 $I_{f0}$ 与短路时产生额定电流所需的励磁电流 $I_{fk}$ 之比：
\begin{equation}
K_c=\frac{I_{f0}}{I_{fk}}\approx\frac{1}{X_d^{*}} .
\end{equation}
短路比大 $\Rightarrow$ $X_d$ 小 $\Rightarrow$ 电压调整率小、并列运行稳定性好，但气隙大、材料多、造价高。汽轮发电机 $K_c\approx0.5\sim0.7$，水轮发电机 $K_c\approx1.0\sim1.4$。

\begin{center}
\begin{tabular}{@{}p{2.2cm}p{3.6cm}p{4.2cm}p{3.0cm}@{}}
\toprule
特性/参数 & 条件 & 形状或关系 & 用途 \\
\midrule
空载特性 & $I=0$，$n=n_1$ & $E_0=f(I_f)$，饱和曲线 & 求 $X_d$、判断饱和程度 \\
短路特性 & $U=0$，$n=n_1$ & $I_k=f(I_f)$，直线 & 求 $X_d$ 与短路比 \\
外特性 & $I_f$、$\cos\varphi$ 恒定 & $U=f(I)$，感性下降 & 电压调整 \\
调整特性 & $U=U_N$、$\cos\varphi$ 恒定 & $I_f=f(I)$ & 励磁调节依据 \\
短路比 & — & $K_c\approx1/X_d^{*}$ & 稳定性与经济性权衡 \\
\bottomrule
\end{tabular}
\end{center}

\corepoint{本章最容易混的两组概念：① \textbf{$X_t$ 与 $X_d$}——隐极机用 $X_t$，凸极机用 $X_d$、$X_q$，两者不能互换代入；② \textbf{空载电动势 $E_0$ 与端电压 $U$}——$E_0$ 是励磁决定的"内部电动势"，$U$ 是带负载后的端电压，两者之差就是电枢反应与漏抗压降。所有同步电机的计算题都围绕 $\dot E_0=\dot U+\mathrm{j}\dot I X$（或其凸极形式）与 $P_{em}=\frac{mE_0U}{X}\sin\delta$ 这两组式子转。}
